\documentclass[11pt]{article}
\usepackage[utf8]{inputenc}
\usepackage{amsmath, amsfonts, amssymb}
\usepackage[margin=2.5cm]{geometry}

% --- MEJORAS DE ESPACIADO ---
\usepackage{parskip}           % Espacio automático entre párrafos sin sangría
\usepackage{setspace}          % Control de interlineado
\setstretch{1}              % Interlineado cómodo de 1.15
\setlength{\jot}{8pt}          % Espaciado vertical interno en fórmulas y alineaciones

% --- PAQUETES PARA CAJAS DE COLOR ---
\usepackage[most]{tcolorbox}
\usepackage{xcolor}

\definecolor{mypurple}{HTML}{800055}      % Violeta oscuro
\definecolor{mylightpurple}{HTML}{FBF0F6} % Violeta claro (fondo)

% Entorno personalizado con aire arriba, abajo e internamente
\newtcolorbox{mynote}[1]{
  colframe=mypurple,
  colback=mylightpurple,
  coltitle=mypurple,
  fonttitle=\bfseries\itshape,
  title=#1,
  sharp corners,
  boxrule=0.8pt,
  titlerule=0.5pt,
  before skip=12pt,
  after skip=12pt,
  boxsep=6pt,
  top=6pt,
  bottom=6pt
}

\title{\textbf{Apuntes de Álgebra I}}
\author{Ivan Ferreyra}
\date{\today}

\begin{document}

\maketitle

\vspace{0.5cm}

\section{Ejercicio guía 8}

\textbf{8.} Sean $a, b \in \mathbb{R}$. Probar que para todo $n \in \mathbb{N}$,

\[
a^n - b^n = (a-b)\sum_{i=1}^{n} a^{i-1}b^{n-i}
\]

Deducir de la fórmula anterior la serie geométrica: para todo $a \neq 1$,

\[
\sum_{i=1}^{n} a^{i-1} = \frac{a^{n}-1}{a-1}
\]

\begin{mynote}{Observación importante}
Acá hay que tener cuidado de no marearse a la hora de leer el ejercicio: nos piden probar la primera ecuación por inducción, y luego la segunda es una deducción directa de la primera.
\end{mynote}

\subsubsection*{Demostración por inducción en $n$}

\textbf{Caso base $p(1)$:}

\[
a^1 - b^1 = (a-b)\sum_{i=1}^{1} a^{i-1}b^{1-i} = (a-b)(a^0 b^0) = (a-b) \cdot 1 = a - b
\]

Por lo tanto, $p(1)$ es verdadero.

\vspace{0.3cm}

\textbf{Paso inductivo:} Asumo que $p(k)$ es verdadero (\textit{Hipótesis Inductiva}):

\[
p(k): a^k - b^k = (a-b)\sum_{i=1}^{k} a^{i-1}b^{k-i}
\]

Quiero probar que $p(k+1)$ es verdadero:

\[
p(k+1): a^{k+1} - b^{k+1} = (a-b)\sum_{i=1}^{k+1} a^{i-1}b^{k+1-i}
\]

Arranco manipulando el lado derecho de la ecuación:

\[
(a-b)\sum_{i=1}^{k+1} a^{i-1}b^{k+1-i} = (a-b) \left( \sum_{i=1}^{k} a^{i-1}b^{k+1-i} + a^{(k+1)-1}b^{k+1-(k+1)} \right)
\]

\begin{mynote}{Estrategia de la sumatoria}
Acá hay que separar el último término de la sumatoria (cuando $i = k+1$) y factorizar la $b$ sobrante para poder hacer aparecer la expresión de la hipótesis inductiva $p(k)$.
\end{mynote}

Separando $b^{k+1-i} = b \cdot b^{k-i}$ y simplificando el último término ($a^k b^0 = a^k$):

\[
= (a-b) \left( b \sum_{i=1}^{k} a^{i-1}b^{k-i} + a^k \right)
\]

Distribuyendo $(a-b)$:

\[
= b \cdot \left[ (a-b)\sum_{i=1}^{k} a^{i-1}b^{k-i} \right] + (a-b)a^k
\]

Aplicando la hipótesis inductiva $p(k)$:

\[
= b(a^k - b^k) + a^{k+1} - ba^k
\]

\[
= ba^k - b^{k+1} + a^{k+1} - ba^k = a^{k+1} - b^{k+1}
\]

Queda demostrado el paso inductivo.

\vspace{0.5cm}

\subsubsection*{Deducción de la serie geométrica}

Tomando la fórmula original $a^n - b^n = (a-b)\sum_{i=1}^{n} a^{i-1}b^{n-i}$ y haciendo $b = 1$:

\[
a^n - 1^n = (a-1)\sum_{i=1}^{n} a^{i-1}(1)^{n-i} \implies a^n - 1 = (a-1)\sum_{i=1}^{n} a^{i-1}
\]

Como $a \neq 1$, podemos dividir por $(a-1)$:

\[
\sum_{i=1}^{n} a^{i-1} = \frac{a^n - 1}{a - 1}
\]

\end{document}